\documentclass[10pt,a4paper]{amsart}
\usepackage[margin=19mm]{geometry}
\usepackage{amsmath,amssymb,mathtools}
\usepackage[scr=rsfs,cal=euler]{mathalfa}
\usepackage{xfrac}
\usepackage[unicode,hidelinks,bookmarks=false]{hyperref}
\newcommand{\ie}{i.e.\ }
\newcommand{\cf}{cf.\ }
\newcommand{\resp}{respectively\ }
\newcommand{\Subsection}{\subsection}
\providecommand{\nobreakdash}{\nobreak\hbox{-}}
\setlength{\emergencystretch}{2em}
\multlinegap0pt
\usepackage{bm}
\usepackage{amssymb}
\usepackage{mathtools}
\usepackage{xcolor}
\newtheorem{theorem}{Theorem}
\title{Inductive Volume Measure: A Geometric Alternative to Hausdorff Measure}
\author{Luis A. Cede\~no-P\'erez}
\date{Source: arXiv:2507.06540v4 (CC BY 4.0)}
\begin{document}
\maketitle

\noindent\emph{Source and licence.} Inductive Volume Measure: A Geometric Alternative to Hausdorff Measure. Version arXiv:2507.06540v4 (CC BY 4.0), distributed by arXiv under the Creative Commons Attribution 4.0 International licence, http://creativecommons.org/licenses/by/4.0/, as recorded in the arXiv licence metadata for this exact version and shown on the abstract page https://arxiv.org/abs/2507.06540v4. Full licence notice: \texttt{licenses/arxiv-2507.06540-CC-BY-4.0.txt}.\par\smallskip

\noindent\textit{Editorial scope.} Counted unit: the article's theorem EqCoarea (source0.tex lines 534--543). The article's complete original proof of it is delivered with it (source0.tex lines 544--568, one \texttt{proof} environment of 25 lines). No mathematics is written, repaired or re-derived: the statement and the proof are the article's own lines, in the article's own notation, and no retained line is substituted or re-flowed.  No further source-local context is needed: every cross-reference of every retained line resolves inside the two delivered spans.  Nothing was omitted from any retained span, so the package carries no omission record.  No retained line contains a citation, so the wrapper carries no bibliography and no bibliography entry.  No retained line contains a figure environment, an image-inclusion command or drawing code, so no image is delivered and the package declares no asset dependency.  Editorial formatting only, in addition: an AMS \texttt{amsart} preamble that restates the article's own declarations and macros used by the retained lines, namely the environments \texttt{theorem} on separate counters, together with the standard packages those lines need.  The retained environments print in this document as Theorem 1 in order of appearance, whereas the article prints the same items with the numbering of its own full document.  The complete native source of the article as distributed in the arXiv e-print for this exact version is bundled unchanged as \texttt{sources/181269/source0.tex}.
\begin{theorem}[Smooth Coarea Formula]
Let $H\colon \mathbb{R}^{n} \to \mathbb{R}$ be a smooth function and for every regular value $t$ let $dV_{t}$ be the volume form of $H^{-1}(\{t\})$. If $dV$ is the volume form of $\mathbb{R}^{n}$ then
\begin{equation*}
    dV = \frac{1}{|\nabla H|} dt\wedge dV_{t}
\end{equation*}
for every regular value $t$. In consequence, for any continuous function $f$ we have that
\begin{equation}\label{EqCoarea}
    \int_{H^{-1}([a,b])} f\;dV = \int_{a}^{b}\int_{H^{-1}(\{t\})} \frac{f}{|\nabla H|}\;dV_{t}\;dt.
\end{equation}
\end{theorem}

\begin{proof}
Let $\Phi_{t}$ be the flow of $\nabla H$ in $H^{-1}(\{t\})$. For sufficiently small $t$ and $x\in H^{-1}(\{t\})$ the flow $\Phi_{t}(x)$ defines smooth coordinates for $\mathbb{R}^{n}$. The one form $dt$ is defined as the unique one form such that
\begin{equation*}
    dt(\nabla H) = |\nabla H|^{2}
\end{equation*}
and
\begin{equation*}
    dt(v) = 0
\end{equation*}
if $v\in T_{x}H^{-1}(\{t\})$. If $\{v_{i}\}_{i=1}^{n-1}$ is an orthonormal frame for $T_{x}H^{-1}(\{t\})$ then $\left\{\frac{\nabla H}{|\nabla H|}\right\}\cup \{v_{i}\}_{i=1}^{n-1}$ is an orthonormal frame for $T_{x}\mathbb{R}^{n}$. It follows that
\begin{align*}
    \frac{1}{|\nabla H|} dt\wedge dV_{t} \left(\frac{\nabla H}{|\nabla H|},v_{1},\ldots,v_{n-1}\right) &= \frac{1}{|\nabla H|}dt\left(\frac{\nabla H}{|\nabla H|}\right) dV_{t} (v_{1},\ldots,v_{n-1})\\
    &= \frac{1}{|\nabla H|^{2}}dt(\nabla H)\\
    &= 1,
\end{align*}
hence
\begin{equation}\label{EqCoareaSmooth}
    dV = \frac{1}{|\nabla H|} dt\wedge dV_{t},
\end{equation}
as stated. For the equality of integrals, note that by Sard's Theorem the set of critical values has measure zero, hence the function
\begin{equation*}
    t \longmapsto \int \frac{f}{|\nabla H|}\;dV_{t}
\end{equation*}
is defined except for a measure zero set, hence it can be integrated and equation (\ref{EqCoareaSmooth}) yields (\ref{EqCoarea}).
\end{proof}

\end{document}
